\section{Method}

% \begin{table}[]
% \centering
% \caption{Mutual Information and Cosine Similarity of CLS token with Patch token mean ($\mu_P$) and Register token mean ($\mu_R$)}
% \begin{tabular}{|c|cc|cc|}
% \hline
% \multirow{2}{*}{Comparison} & \multicolumn{2}{c|}{Mutual Information} & \multicolumn{2}{c|}{Cosine Similarity} \\ \cline{2-5} 
%                             & ID                 & OOD                & ID                 & OOD               \\
% CLS vs $\mu_P$              & 0.256              & 0.286              & 0.547              & 0.550             \\
% CLS vs $\mu_R$              & 0.018              & 0.004              & 0.072              & 0.038             \\ \hline
% \end{tabular}

% \label{tab:mutual_info}
% \end{table}

\noindent\textbf{Preliminaries:}
Vision Transformers (ViTs) process images by dividing them into patches and viewing them as a sequence of tokens, where each token corresponds to a patch in the image. An additional token, the \texttt{[CLS]} token, is appended to the patch sequence to aggregate the global information from the image patches. The \texttt{[CLS]} embedding is often used as the global representation of the image in downstream tasks. %Our proposed method leverages the auxiliary information in register tokens to augment the information found in the \texttt{[CLS]} token, enhancing out-of-distribution (OOD) detection performance. The approach is designed to be simple yet effective, with no additional computational overhead.

Let $X \in \mathbb{R}^{H \times W \times 3}$ represent a $3$ channel RGB input image, where $H$ and $W$ are the height and width of the image, respectively. A ViT backbone, denoted by $F$, divides $X$ into $L$ patches, each of size $K \times K$, and projects them into $D$-dimensional embeddings. The \texttt{[CLS]} token $c \in \mathbb{R}^D$ is added to aggregate global information. The backbone produces:  

\begin{itemize}
    \item \texttt{[CLS]} token $c \in \mathbb{R}^D$
    \item Patch tokens $P = \{p_1 \ldots p_L \} \in \mathbb{R}^{L \times D} $
    \item Register tokens $R = \{r_1 \ldots r_M \} \in \mathbb{R}^{M \times D}$
\end{itemize}

where $D$ is the hidden dimensionality of the transformer, and $L$ and $M$ are the numbers of patch tokens and register tokens, respectively. 

In typical ViT models, downstream tasks use either the \texttt{[CLS]} token or the mean of patch tokens $\mu_P = \frac{1}{L} \sum_{j=1}^L p_j$ \cite{Zhai_2022_CVPR,li2024inverse}. In Dino-v2, the \texttt{[CLS]} token is concatenated with $\mu_p$, and fed to a linear classifier $h_{\theta}$ which is optimized with the following objective 

\begin{align}
\min_{\theta} \frac{1}{N} \sum_{i=1}^{N} \ell(h_{\theta}(f_i), y_i)
\label{eq:ERM}
\end{align}

where $N$ is the total number of samples, $f_i = [c_i; \mu_P^i]$ is the concatenated  \texttt{[CLS]} token \( c_i \) and the mean of the patch tokens \( \mu_P^i \) for the \(i^\text{th}\) sample, $h_{\theta}(f_i) = f_i^\top \theta$ denotes the linear probe with parameters \( \theta \), and \( \ell(\cdot, \cdot) \) is the cross-entropy loss computed between the prediction \( h_{\theta}(f_i) \) and the target \( y_i \). \\


\noindent\textbf{Proposed Approach:}
We propose to leverage the auxiliary information in register tokens to augment the information found in the \texttt{[CLS]} token, thereby enhancing out-of-distribution (OOD) detection performance. As previously stated, Darcet \textit{et al.} \cite{darcet2024vision} noted that register tokens capture global information similar to the \texttt{[CLS]} token, by training linear probes on register tokens without observing significant drop in ID performance. Furthermore, Zhang et al \cite{zhang2023learning} showed that although similar representations from models trained with different random seeds yield comparable ID performance, they can encode different information under OOD conditions. Crucially, combining these representations leads to richer feature sets that outperform either representation alone in OOD tasks.

\renewcommand{\arraystretch}{1.4}
\begin{table*}[ht]
\centering
\caption{\textbf{Evaluating ID and OOD generalization as well anomaly rejection performance across different ViT architectures.} We highlight the best performing method in every case with green. While the ID accuracies remain comparable to a baseline variant, we find that our proposed approach consistently outperforms the existing baselines on OOD generalization as well as anomaly rejection across a variety of benchmarks and architectures.}

\renewcommand{\arraystretch}{1.4}
\resizebox{1\linewidth}{!}{
\begin{tabular}{|cccccccccccccc|}
\hline
\rowcolor[gray]{0.92} % Light grey background
\multicolumn{14}{|c|}{ViT-Giant}                                                                                                                                                                                                                              \\ \hline
\multicolumn{1}{|c|}{\multirow{2}{*}{Method}} & \multicolumn{1}{c|}{\multirow{2}{*}{ \shortstack{Dino-v2 \\ with registers}}} & \multicolumn{1}{c|}{\multirow{2}{*}{ID Acc}} & \multicolumn{3}{c|}{ImageNet OOD}                                     & \multicolumn{8}{c|}{Anomaly Rejection (FPR $\downarrow$ / AUROC $\uparrow$)}                                                                          \\ \cline{4-14} 
\multicolumn{1}{|c|}{}                        & \multicolumn{1}{c|}{}                           & \multicolumn{1}{c|}{}                         & In-A           & In-R           & \multicolumn{1}{c|}{In-S}           & \multicolumn{1}{c|}{Score} & DTD         & SVHN        & Places-365  & LSUN        & LSUN-R      & ISUN        & Mean        \\ \hline
\multicolumn{1}{|c|}{\texttt{CLS} ; $\mu_P$}       & \multicolumn{1}{c|}{\ding{53}}                          & \multicolumn{1}{c|}{86.5}                    & 73.38          & 77.63           & \multicolumn{1}{c|}{61.24}          & \multicolumn{1}{c|}{\multirow{3}{*}{MSP}}   & 48.81/85.32 & 22.68/95.71   &  { {\color[HTML]{0f9845}\textbf{59.65}}}/83.0  & 18.76/96.04 & 43.32/90.95  & 41.9/90.32 & 39.19/90.22 \\
\multicolumn{1}{|c|}{\texttt{CLS} ; $\mu_P$}  & \multicolumn{1}{c|}{\color[HTML]{237A1E}\faCheckSquare{ }}                          & \multicolumn{1}{c|}{\color[HTML]{0f9845}\textbf{87.1}}                    & 75.38          & 79.05          & \multicolumn{1}{c|}{62.90}           & \multicolumn{1}{c|}{}    & 46.42/86.16  & 27.16/94.57 & 58.58/{{\color[HTML]{0f9845}\textbf{83.09}}} & 16.3/96.5 & 40.41/91.14  & 38.27/91.25 & 37.86/90.45 \\
\multicolumn{1}{|c|}{\texttt{CLS} ; $\mu_R$}  & \multicolumn{1}{c|}{\color[HTML]{237A1E}\faCheckSquare{ } }                          & \multicolumn{1}{c|}{86.57}   &  {\color[HTML]{0f9845}\textbf{78.09}} &  {\color[HTML]{0f9845}\textbf{80.51}} & \multicolumn{1}{c|}{ {\color[HTML]{0f9845}\textbf{64.43}}} & \multicolumn{1}{c|}{}    &  {\color[HTML]{0f9845}\textbf{45.53}}/ {\color[HTML]{0f9845}\textbf{86.07}} &  {\color[HTML]{0f9845}\textbf{16.71}}/ {\color[HTML]{0f9845}\textbf{96.71}} & {58.88}/82.79 &  {\color[HTML]{0f9845}\textbf{11.78}}/{\color[HTML]{0f9845}\textbf{97.48}} &  {\color[HTML]{0f9845}\textbf{37.06}}/{\color[HTML]{0f9845}\textbf{91.59}} &  {\color[HTML]{0f9845}\textbf{33.76}}/{\color[HTML]{0f9845}\textbf{92.22}} &  {\color[HTML]{0f9845}\textbf{33.95}}/{\color[HTML]{0f9845}\textbf{91.14}} \\\hline
\multicolumn{1}{|c|}{\texttt{CLS} ; $\mu_P$}       & \multicolumn{1}{c|}{\ding{53}}                          & \multicolumn{1}{c|}{-}                        & -              & -              & \multicolumn{1}{c|}{-}              & \multicolumn{1}{c|}{\multirow{3}{*}{Energy}} & 31.83/{{\color[HTML]{0f9845}\textbf{92.12}}} & 10.11/97.66 & 43.04/{{\color[HTML]{0f9845}\textbf{89.72}}} & 6.34/98.60  & 30.30/94.92 & 26.73/94.74 & 24.73/94.63 \\ 
\multicolumn{1}{|c|}{\texttt{CLS} ; $\mu_P$}  & \multicolumn{1}{c|}{\color[HTML]{237A1E}\faCheckSquare{ }}                          & \multicolumn{1}{c|}{-}                        & -              & -              & \multicolumn{1}{c|}{-}              & \multicolumn{1}{c|}{} & {{\color[HTML]{0f9845}\textbf{30.48}}}/91.45 & 6.81/98.52 & 43.17/88.84 & 5.03/98.85  & 21.77/{\color[HTML]{0f9845}\textbf{95.78}} & 20.25/95.70 & 21.25/{\color[HTML]{0f9845}\textbf{94.86}}\\ 
\multicolumn{1}{|c|}{\texttt{CLS} ; $\mu_R$}  & \multicolumn{1}{c|}{\color[HTML]{237A1E}\faCheckSquare{ }}                          & \multicolumn{1}{c|}{-}                        & -              & -              & \multicolumn{1}{c|}{-}              & \multicolumn{1}{c|}{} & 32.41/90.64 & {\color[HTML]{0f9845}\textbf{4.96}}/{\color[HTML]{0f9845}\textbf{98.92}}  & {\color[HTML]{0f9845}\textbf{45.5}}/{\color[HTML]{0f9845}\textbf{87.64}}  & {\color[HTML]{0f9845}\textbf{3.85}}/{\color[HTML]{0f9845}\textbf{99.06}}  & {\color[HTML]{0f9845}\textbf{19.95}}/95.76 & {\color[HTML]{0f9845}\textbf{18.48}}/{\color[HTML]{0f9845}\textbf{95.95}} & {\color[HTML]{0f9845}\textbf{20.86}}/94.66\\ \hline
\end{tabular}
}

\vspace{0.5em} 
\renewcommand{\arraystretch}{1.4}
\resizebox{1\linewidth}{!}{
\begin{tabular}{|cccccccccccccc|}
\hline
\rowcolor[gray]{0.92}
\multicolumn{14}{|c|}{ViT-Large}                                                                                                                                                                                                                     \\ \hline
\multicolumn{1}{|c|}{\multirow{2}{*}{Method}} & \multicolumn{1}{c|}{\multirow{2}{*}{ \shortstack{Dino-v2 \\ with registers}} } & \multicolumn{1}{c|}{\multirow{2}{*}{ID Acc}} & \multicolumn{3}{c|}{ImageNet OOD}                                    & \multicolumn{8}{c|}{Anomaly Rejection (FPR $\downarrow$ / AUROC $\uparrow$)}                                                                          \\ \cline{4-14} 
\multicolumn{1}{|c|}{}                        & \multicolumn{1}{c|}{}                           & \multicolumn{1}{c|}{}                             & In-A           & In-R          & \multicolumn{1}{c|}{In-S}           & \multicolumn{1}{c|}{Score} & DTD         & SVHN        & Places-365  & LSUN        & LSUN-R      & ISUN        & Mean        \\ \hline
\multicolumn{1}{|c|}{\texttt{CLS} ; $\mu_P$}                & \multicolumn{1}{c|}{\ding{53}}                          & \multicolumn{1}{c|}{86.3}                            & 71.24          & 74.52          & \multicolumn{1}{c|}{59.46}          & \multicolumn{1}{c|}{\multirow{3}{*}{MSP}} & 51.4/84.59 & 27.2/94.6 & 59.54/83.24 & 20.63/95.49 & 40.81/91.61 & 40.46/90.9    & 40.1/90.07 \\
\multicolumn{1}{|c|}{\texttt{CLS} ; $\mu_P$}           & \multicolumn{1}{c|}{\color[HTML]{237A1E}\faCheckSquare{ }}                          & \multicolumn{1}{c|}{\color[HTML]{0f9845}\textbf{86.7}}                            & 73.65          & 76.31         & \multicolumn{1}{c|}{61.86}          & \multicolumn{1}{c|}{}                     &48.12/85.67 & 22.44/95.42  & 58.13/83.54 & 17.36/96.32 & 37.35/92.39 & 37.49/91.48 & 36.82/90.81  \\
\multicolumn{1}{|c|}{\texttt{CLS} ; $\mu_R$}           & \multicolumn{1}{c|}{\color[HTML]{237A1E}\faCheckSquare{ }}                          & \multicolumn{1}{c|} {85.82}                            & {\color[HTML]{0f9845}\textbf{75.33}} & {\color[HTML]{0f9845}\textbf{79.2}} & \multicolumn{1}{c|}{\color[HTML]{0f9845}{\textbf{62.82}}} & \multicolumn{1}{c|}{}                     & {\color[HTML]{0f9845}\textbf{43.71}}/{\color[HTML]{0f9845}\textbf{86.25}} & {\color[HTML]{0f9845}\textbf{15.93}}/{\color[HTML]{0f9845}\textbf{96.76}} & {\color[HTML]{0f9845}\textbf{55.88}}/{\color[HTML]{0f9845}\textbf{83.35}} & {\color[HTML]{0f9845}\textbf{12.65}}/{\color[HTML]{0f9845}\textbf{97.12}} & {\color[HTML]{0f9845}\textbf{33.49}}/{\color[HTML]{0f9845}\textbf{92.38}} & {\color[HTML]{0f9845}\textbf{31.35}}/{\color[HTML]{0f9845}\textbf{92.77}} & {\color[HTML]{0f9845}\textbf{32.17}}/{\color[HTML]{0f9845}\textbf{91.44}}
 \\ \hline
\multicolumn{1}{|c|}{\texttt{CLS} ; $\mu_P$}                & \multicolumn{1}{c|}{\ding{53}}                          & \multicolumn{1}{c|}{-}                            & \textbf{-}     & -             & \multicolumn{1}{c|}{-}              & \multicolumn{1}{c|}{\multirow{3}{*}{Energy}} & 47.27/87.26 & 48.95/93.03 & 55.77/86.15 & 17.8/96.13 & 49.99/91.78 & 42.85/91.91 & 43.77/91.04  \\ 
\multicolumn{1}{|c|}{\texttt{CLS} ; $\mu_P$}           & \multicolumn{1}{c|}{\color[HTML]{237A1E}\faCheckSquare{ }}                          & \multicolumn{1}{c|}{-}                            & -              & -             & \multicolumn{1}{c|}{-}              & \multicolumn{1}{c|}{}   & 35.83/90.51  & 10.35/97.73 & 45.21/{\color[HTML]{0f9845}\textbf{89.05}} & 8.08/98.18 & 24.83/95.76 & 24.56/95.21 & 24.81/94.41 \\ 
\multicolumn{1}{|c|}{\texttt{CLS} ; $\mu_R$}           & \multicolumn{1}{c|}{\color[HTML]{237A1E}\faCheckSquare{ }}                          & \multicolumn{1}{c|}{-}                            & -              & -             & \multicolumn{1}{c|}{-}              & \multicolumn{1}{c|}{}   & {\color[HTML]{0f9845}\textbf{34.06}}/{\color[HTML]{0f9845}\textbf{90.87}} & {\color[HTML]{0f9845}\textbf{7.05}}/{\color[HTML]{0f9845}\textbf{98.55}} & {\color[HTML]{0f9845}\textbf{43.44}}/88.79 & {\color[HTML]{0f9845}\textbf{6.86}}/{\color[HTML]{0f9845}\textbf{98.5}} & {\color[HTML]{0f9845}\textbf{19.24}}/{\color[HTML]{0f9845}\textbf{96.21}} & {\color[HTML]{0f9845}\textbf{18.87}}/{\color[HTML]{0f9845}\textbf{96.15}} & {\color[HTML]{0f9845}\textbf{21.59}}/{\color[HTML]{0f9845}\textbf{94.84}}
\\ \hline
\end{tabular}

}


\vspace{0.5em}
\renewcommand{\arraystretch}{1.4}
\resizebox{1\linewidth}{!}{
\begin{tabular}{|cccccccccccccc|}
\hline
\rowcolor[gray]{0.92}
\multicolumn{14}{|c|}{ViT-Base}                                                                                                                                                                                                                       \\ \hline
\multicolumn{1}{|c|}{\multirow{2}{*}{Method}} & \multicolumn{1}{c|}{\multirow{2}{*}{ \shortstack{Dino-v2 \\ with registers}}} & \multicolumn{1}{c|}{\multirow{2}{*}{ID Acc}} & \multicolumn{3}{c|}{ImageNet OOD}                                     & \multicolumn{8}{c|}{Anomaly Rejection (FPR $\downarrow$ / AUROC $\uparrow$)}                                                                          \\ \cline{4-14} 
\multicolumn{1}{|c|}{}                        & \multicolumn{1}{c|}{}                           & \multicolumn{1}{c|}{}                             & In-A           & In-R           & \multicolumn{1}{c|}{In-S}           & \multicolumn{1}{c|}{Score} & DTD         & SVHN        & Places-365  & LSUN        & LSUN-R      & ISUN        & Mean        \\ \hline
\multicolumn{1}{|c|}{\texttt{CLS} ; $\mu_P$}                & \multicolumn{1}{c|}{\ding{53}}                          & \multicolumn{1}{c|}{\color[HTML]{0f9845}\textbf{84.5} }                           & 55.54 & 63.9         & \multicolumn{1}{c|}{50.92}          & \multicolumn{1}{c|}{\multirow{3}{*}{MSP}}   & 55.66/83.95 & {\color[HTML]{0f9845}\textbf{18.35}}/{\color[HTML]{0f9845}\textbf{96.44}}  & 64.65/81.76 & 21.94/95.42 & 39.51/91.57 & 39.9/91.24 & 40.01/90.06 \\
\multicolumn{1}{|c|}{\texttt{CLS} ; $\mu_P$}           & \multicolumn{1}{c|}{\color[HTML]{237A1E}\faCheckSquare{ }}                          & \multicolumn{1}{c|}{84.21}                           & 54.56          & 65.02          & \multicolumn{1}{c|}{53.61}          & \multicolumn{1}{c|}{}    & 53.12/84.45  & 20.31/96.27 & 62.87/82.34 & 16.84/96.31 & 33.83/92.98 & 35.38/92.24 & 37.06/90.76  \\
\multicolumn{1}{|c|}{\texttt{CLS} ; $\mu_R$}           & \multicolumn{1}{c|}{\color[HTML]{237A1E}\faCheckSquare{ }}                          & \multicolumn{1}{c|}{83.84}                            & {\color[HTML]{0f9845}\textbf{57.87}}          &{\color[HTML]{0f9845}\textbf{68.75}} & \multicolumn{1}{c|}{\color[HTML]{0f9845}\textbf{55.14}} & \multicolumn{1}{c|}{}    & {\color[HTML]{0f9845}\textbf{48.09}}/{\color[HTML]{0f9845}\textbf{85.42}}& 19.83/95.8 & {\color[HTML]{0f9845}\textbf{59.6}}/{\color[HTML]{0f9845}\textbf{82.48}} & {\color[HTML]{0f9845}\textbf{13.55}}/{\color[HTML]{0f9845}\textbf{97.09}} & {\color[HTML]{0f9845}\textbf{30.24}}/{\color[HTML]{0f9845}\textbf{93.66}} & {\color[HTML]{0f9845}\textbf{29.56}}/{\color[HTML]{0f9845}\textbf{93.36}} & {\color[HTML]{0f9845}\textbf{33.48}}/{\color[HTML]{0f9845}\textbf{91.3}}
 \\\hline
\multicolumn{1}{|c|}{\texttt{CLS} ; $\mu_P$}                & \multicolumn{1}{c|}{\ding{53}}                          & \multicolumn{1}{c|}{-}                            & \textbf{-}     & -              & \multicolumn{1}{c|}{-}              & \multicolumn{1}{c|}{\multirow{3}{*}{Energy}} & 60.09/85.47 & 47.68/93.64 & 74.61/82.53 & 33.55/94.50 & 70.7/88.61 & 64.44/89.28 & 58.51/89.01 \\
\multicolumn{1}{|c|}{\texttt{CLS} ; $\mu_P$}           & \multicolumn{1}{c|}{\color[HTML]{237A1E}\faCheckSquare{ }}                          & \multicolumn{1}{c|}{-}                            & \textbf{-}     & \textbf{-}     & \multicolumn{1}{c|}{\textbf{-}}     & \multicolumn{1}{c|}{} & 42.32/89.58 & 18.22/96.34 & 52.38/87.91 & 8.73/98.06     & 24.85/95.82 & 24.03/95.44 & 28.59/93.86 \\
\multicolumn{1}{|c|}{\texttt{CLS} ; $\mu_R$}           & \multicolumn{1}{c|}{\color[HTML]{237A1E}\faCheckSquare{ }}                          & \multicolumn{1}{c|}{-}                            & -              & -              & \multicolumn{1}{c|}{-}              & \multicolumn{1}{c|}{} & {\color[HTML]{0f9845}\textbf{37.54}}/{\color[HTML]{0f9845}\textbf{90.71}} & {\color[HTML]{0f9845}\textbf{12.02}}/{\color[HTML]{0f9845}\textbf{97.46}} & {\color[HTML]{0f9845}\textbf{46.38}}/{\color[HTML]{0f9845}\textbf{88.83}} & {\color[HTML]{0f9845}\textbf{7.07}}/{\color[HTML]{0f9845}\textbf{98.57}} & {\color[HTML]{0f9845}\textbf{17.91}}/{\color[HTML]{0f9845}\textbf{96.97}} & {\color[HTML]{0f9845}\textbf{16.87}}/{\color[HTML]{0f9845}\textbf{96.74}} & {\color[HTML]{0f9845}\textbf{22.96}}/{\color[HTML]{0f9845}\textbf{94.88}}
 \\ \hline
\end{tabular}
}
\label{tab:results}
\end{table*}



Building on this insight, we concatenate the \texttt{[CLS]} token with the mean of register tokens, $\mu_R = \frac{1}{M} \sum_{k=1}^M r_k$, as opposed to discarding the register tokens as described in previous works  \cite{darcet2024vision, hoptimus0}.  This combined representation exploits the auxiliary information captured by the register tokens, allowing us to create a richer feature representation. The concatenated representation is then fed to the linear classifier $h_{\theta}$ which is optimized with the same objective as presented in Eq. \ref{eq:ERM} except that the $f_i$ is now given by  $f_i = [c_i; \mu_R^i]$. Figure. \ref{fig:method} provides an overview of the proposed approach. 

% \begin{align*}
%     Output = W[c; \mu_R]
% \end{align*}



%Conventional approaches use either the CLS token $c$, the mean of patch tokens $\mu_P = \frac{1}{N} \sum_{i=1}^N p_i$ \cite{Zhai_2022_CVPR,li2024inverse} , or a combination of both \cite{oquab2023dinov2} for downstream tasks. In contrast, our method utilizes both the CLS token and the mean of register tokens $\mu_R = \frac{1}{M} \sum_{i=1}^M r_i$, the latter of which is typically discarded in previous works \cite{darcet2024vision, hoptimus0}.

% To verify the complementarity of these tokens, we conducted an analysis using 1500 random images from ImageNet \cite{russakovsky2015imagenet} for ID data, and a mixture from ImageNet-R \cite{hendrycks2021many} and ImageNet-S \cite{wang2019learning} for OOD data. We computed two metrics: \textbf{Mutual Information} (MI) \footnote{Computed using scikit-learn mutual\_info\_regression} and \textbf{Cosine Similarity} (CS). These metrics were calculated  between the CLS token $c$ and $\mu_R$, as well as between $c$ and $\mu_P$. As shown in Table \ref{tab:mutual_info}, our findings reveal low MI and CS between $c$ and $\mu_R$ (MI of 0.018 for ID), but high MI and CS between $c$ and $\mu_P$ (MI of 0.256 for ID). Importantly, this pattern persists in both ID and OOD scenarios. This empirical evidence supports our hypothesis that Register tokens contain information complementary to the CLS token, suggesting that their combination could yield a richer feature representation for downstream tasks.


Note our method does not introduce additional computational overhead because the ViT backbone is already trained and kept frozen during the experiments. We only need to train a linear layer, however, just like the baseline approach that concatenates the \texttt{[CLS]} token with the patch token mean. Since both methods involve concatenating two vectors ($[\texttt{CLS}; \mu_R]$ vs $[\texttt{CLS}; \mu_P]$), the linear layer remains the same size. Thus, given a backbone already trained with register tokens, the only step required is training this linear layer, making our approach computationally efficient.